\section{Actual original power comparisons on the unit open}

Fix the original test configuration, an original line $M$ on its total
space, a line $L$ on the original fibre, and the given whole action and
fibre-one comparisons $\alpha,e_1$.  All pullbacks below use these same
original maps.  Write $\beta(M,L;\alpha,e_1)$ for the resulting whole
unit-open line isomorphism.

\begin{lemma}[The constructed original power data give the actual power isomorphism]
\label{lem:native-tc-nef-fs-unit-open-power-coherence}
\leanok
\uses{lem:native-tc-nef-fs-whole-tensor-comparisons, thm:native-test-configuration-line-trivialization, prop:native-test-configuration-line-powers, thm:actual-power-pullback-composition, lem:actual-power-pullback-naturality}
\lean{EnvelopeProperty.NativeTestConfiguration.TestConfiguration.unitOpenLineIso_power_comparisons}
For every natural $n$, the whole unit-open line isomorphism constructed
from the original power line $M^{\otimes n}$, the original generic power
$L^{\otimes n}$ and the internally constructed power comparisons
$\alpha_n,e_{1,n}$ equals the canonical power of
$\beta(M,L;\alpha,e_1)$, with the actual pullback-power comparisons on
its source and target.
\end{lemma}
\begin{proof}
Expand the whole unit-open isomorphism into its original nine factors:
action-map congruence, inverse action pullback composition, pulled action
comparison, projection pullback composition, projection-map congruence,
inverse fibre pullback composition, pulled fibre-one comparison, original
fibre-map pullback composition and original-projection congruence.
Actual pullback-power coherence gives a commuting square for each map
congruence and each pullback-composition factor.  Pullback-power
naturality gives the square for each pulled original comparison; the
constructed power comparisons are exactly the required conjugates.
Compose these nine commuting squares.  Functoriality of actual tensor
powers identifies their lower composite with the power of the original
whole unit-open isomorphism.  Cancel the target pullback-power
isomorphism to obtain the claimed whole isomorphism identity.
No equality of potentials or of section norms is used.
\end{proof}

\begin{lemma}[The constructed power data on the original product]
\label{lem:native-tc-nef-fs-product-power-coherence}
\leanok
\uses{lem:native-tc-nef-fs-unit-open-power-coherence, def:native-tc-product-original-line-iso}
\lean{EnvelopeProperty.NativeTestConfiguration.TestConfiguration.productTotalLineIso_power_comparisons}
The whole product line isomorphism constructed from the same power data
equals the original canonical product power line isomorphism.
\end{lemma}
\begin{proof}
Insert the preceding whole unit-open identity into the definition of
the original product line isomorphism.  The remaining pullback-composition
and original-projection congruence factors are literally the factors
defining the original canonical product power line isomorphism.
\end{proof}

\begin{remark}
These identities retain the original action, fibre-one isomorphism and
all constants carried by them.  They are structural input to the actual
norm-factor power law; they do not assert that law or any numerical-class
comparison on their own.
\end{remark}
